% 08-parser-stmts.tex — Chapter 8: Parser: Statements and Declarations
\chapter{Parser: Statements and Declarations}
\label{chap:08-parser-stmts}
\index{parser!statements}
\index{parser!declarations}

\pnum
Statement parsing is in \srcref{src/parser/stmt.h}; declaration (top-level) parsing
is in \srcref{src/parser/decl.h}.  Both files use the framework described in
\cref{chap:06-parser-framework}.

% -------------------------------------------------------------------------
\section{Statement list parsing}
\label{sec:par-stmtlist}
% -------------------------------------------------------------------------

\pnum
\cname{parse\_stmt\_list(arena, parser)} (\srcref{src/parser/stmt.h:72})
accumulates statements into a \cname{StmtList} until it encounters \texttt{\}}
or \cname{TOKEN\_EOF}.  After each statement it requires an end-of-line token
(\cname{TOKEN\_EOL}), then skips any subsequent blank lines.

\begin{algorithm}
\textbf{parse\_stmt\_list}(\textit{arena, parser}):
\begin{enumerate}
\item Skip leading EOL tokens.
\item While not at \texttt{\}} or EOF:
  \begin{enumerate}[(a)]
  \item Parse one statement via \cname{parse\_stmt}.
  \item Append to the list.
  \item If the next token is \texttt{\}}, stop.
  \item Require EOL (newline terminates the statement).
  \item Skip further EOL tokens.
  \end{enumerate}
\item Return the list.
\end{enumerate}
\end{algorithm}

% -------------------------------------------------------------------------
\section{Statement dispatch}
\label{sec:par-stmt-dispatch}
% -------------------------------------------------------------------------

\pnum
\cname{parse\_stmt} (\srcref{src/parser/stmt.h}) dispatches on the current token
to the appropriate parse function:

\begin{center}
\begin{tabular}{lll}
\toprule
\textbf{Token} & \textbf{Parse function} & \textbf{Stmt kind} \\
\midrule
\lain{var}      & \cname{parse\_var\_stmt}      & \cname{STMT\_VAR} \\
\lain{if}       & \cname{parse\_if\_stmt}       & \cname{STMT\_IF} \\
\lain{elif}     & (handled inside if)           & \cname{STMT\_IF} (chained) \\
\lain{for}      & \cname{parse\_for\_stmt}      & \cname{STMT\_FOR} \\
\lain{while}    & \cname{parse\_while\_stmt}    & \cname{STMT\_WHILE} \\
\lain{case}     & \cname{parse\_match\_stmt}    & \cname{STMT\_MATCH} \\
\lain{return}   & (inline)                      & \cname{STMT\_RETURN} \\
\lain{unsafe}   & \cname{parse\_unsafe\_stmt}   & \cname{STMT\_UNSAFE} \\
\lain{defer}    & \cname{parse\_defer\_stmt}    & \cname{STMT\_DEFER} \\
\lain{use}      & \cname{parse\_use\_stmt}      & \cname{STMT\_USE} \\
\lain{continue} & \cname{parse\_continue\_stmt} & \cname{STMT\_CONTINUE} \\
\lain{break}    & \cname{parse\_break\_stmt}    & \cname{STMT\_BREAK} \\
\lain{comptime} & \cname{parse\_comptime\_stmt} & \cname{STMT\_COMPTIME\_IF} \\
identifier      & \cname{parse\_assign\_stmt} or \cname{parse\_decl\_stmt} & \cname{STMT\_ASSIGN}/\cname{STMT\_VAR} \\
otherwise       & \cname{parse\_expr\_stmt}     & \cname{STMT\_EXPR} \\
\bottomrule
\end{tabular}
\end{center}

% -------------------------------------------------------------------------
\section{Variable declarations}
\label{sec:par-stmt-var}
\index{var statement}
% -------------------------------------------------------------------------

\pnum
There are two syntactic forms for variable declarations:

\begin{enumerate}
\item \textbf{Explicit \lain{var}}: \lain{var name [Type] = expr}.  The \lain{var}
  keyword is required for mutable bindings.  \cname{parse\_var\_stmt} expects the
  \lain{var} keyword, then parses name, optional type annotation, and required
  initializer.
\item \textbf{Implicit}: \lain{name Type = expr}.  When \cname{parse\_stmt} sees
  an identifier that is not a known keyword, it calls \cname{parse\_decl\_stmt}
  which treats the identifier as the variable name, followed by a type annotation
  and an initializer.
\end{enumerate}

\pnum
Both forms require an initializer; \lain{= undefined} is allowed as a
bypass marker.  Immutable declarations (\lain{name Type = undefined}) are
rejected at parse time.

% -------------------------------------------------------------------------
\section{Assignment statements}
\label{sec:par-stmt-assign}
\index{assignment statements}
% -------------------------------------------------------------------------

\pnum
An assignment is parsed by \cname{parse\_assign\_stmt} when the current token is
an identifier and the lookahead after parsing a target expression is \texttt{=}
(or a compound assignment operator like \texttt{+=}, \texttt{-=}, etc.).  The
target is an arbitrary left-hand side expression (identifier, member access,
index expression).  Compound assignments are desugared by the parser into
\cname{target = target op rhs} before AST construction.

% -------------------------------------------------------------------------
\section{Control-flow statements}
\label{sec:par-stmt-cf}
% -------------------------------------------------------------------------

\subsection{if / elif / else}
\label{sec:par-stmt-if}
\index{if statement}

\pnum
\cname{parse\_if\_stmt} parses the condition expression, the \texttt{\{} block,
then checks for \lain{elif} or \lain{else}.  An \lain{elif} is desugared into a
nested \cname{STMT\_IF} in the else branch.  The else branch is either a
\cname{StmtList} (from \lain{else \{ ... \}}) or the single nested
\cname{STMT\_IF} from \lain{elif}.

\subsection{for}
\label{sec:par-stmt-for}
\index{for statement}

\pnum
\lain{for [idx,] val in iterable \{ body \}} is parsed by \cname{parse\_for\_stmt}.
The optional index name is separated from the value name by a comma.
The iterable is an arbitrary expression.  The body is a statement list.

\subsection{while}
\label{sec:par-stmt-while}
\index{while statement}

\pnum
\cname{parse\_while\_stmt} parses \lain{while cond [decreasing measure] \{ body \}}.
After the condition, the parser checks for the \cname{TOKEN\_KEYWORD\_DECREASING}
token.  If present, the next expression is the termination measure, stored in
\cname{StmtWhile.measure}; otherwise \cname{measure} is \texttt{NULL}.

\pnum
A \texttt{NULL} measure means the while loop is unbounded and is only legal in a
\lain{proc} (procedure) context.  The sema pass enforces this constraint.

\subsection{Match statements}
\label{sec:par-stmt-match}
\index{match statement}

\pnum
\lain{case [&] scrutinee \{ arms \}} is parsed by \cname{parse\_match\_stmt}.
The optional \texttt{\&} marks the match as borrowed (\cname{is\_borrowed = true}).
Each arm is a comma-separated list of patterns followed by \texttt{=>} and a body.
The wildcard arm uses the \lain{else} keyword as the pattern.

% -------------------------------------------------------------------------
\section{Declaration parsing}
\label{sec:par-decls}
\index{parser!declarations}
% -------------------------------------------------------------------------

\pnum
\cname{parse\_module(arena, parser)} in \srcref{src/parser/decl.h} is the top-level
entry point.  It calls \cname{parse\_decl} in a loop until \cname{TOKEN\_EOF}.

\pnum
Top-level declarations parsed by \cname{parse\_decl}:

\begin{center}
\begin{tabular}{lll}
\toprule
\textbf{Keyword} & \textbf{Decl kind} & \textbf{Payload} \\
\midrule
\lain{func}      & \cname{DECL\_FUNCTION}   & \cname{DeclFunction} \\
\lain{proc}      & \cname{DECL\_PROCEDURE}  & \cname{DeclFunction} \\
\lain{struct}    & \cname{DECL\_STRUCT}     & \cname{DeclStruct} \\
\lain{enum}      & \cname{DECL\_ENUM}       & \cname{DeclEnum} \\
\lain{import}    & \cname{DECL\_IMPORT}     & \cname{DeclImport} \\
\lain{extern}    & \cname{DECL\_EXTERN\_*}  & \cname{DeclFunction} \\
\lain{c\_include}& \cname{DECL\_C\_INCLUDE} & \cname{DeclCInclude} \\
\lain{type}      & \cname{DECL\_TYPE\_ALIAS}& \cname{DeclTypeAlias} \\
\bottomrule
\end{tabular}
\end{center}

\pnum
Function and procedure parsing share the same \cname{DeclFunction} payload.
The body parser calls \cname{parse\_stmt\_list} for the block.
Optional \lain{pre}/\lain{post} contract clauses are parsed after the return
type and before the body block.

\pnum
Struct field declarations inside a \texttt{struct} body are parsed as a list of
\cname{DECL\_VARIABLE} nodes.  Each field may carry an optional \lain{in field}
annotation (for bounds-proven indices) and equation-style constraints.

\pnum
Enum declarations parse a list of variants.  Each variant is either a simple
identifier or an identifier followed by a brace-delimited list of fields
(forming an algebraic data type variant).
